\documentclass[journal=jctcce,manuscript=suppinfo]{achemso}
\usepackage{booktabs}
\usepackage{longtable}
\usepackage{siunitx}
\sisetup{scientific-notation=false}
% Filled from the six-run logs on eec66fe3.
% Median is the mean of the 3rd and 4th sorted values.
% IQR is the 5th sorted value minus the 2nd.
% A miss is charged twice the ceiling, 219506.

\newcommand{\HitsSix}{6}

\newcommand{\TempReduced}{0.8}

\newcommand{\CommitSha}{eec66fe30a12664a73bd4c091af1cb7fabb95c47}

\newcommand{\LibCommit}{eec66fe30a12664a73bd4c091af1cb7fabb95c47}

\newcommand{\MainCommit}{46c8c867d23d6dba75a878a7614144687fb31542}

\newcommand{\CallCeiling}{109753}

\newcommand{\MissPar}{219506}

\newcommand{\HitLjxxxviii}{4}

\newcommand{\MedLjxxxviii}{39147.5}

\newcommand{\IqrLjxxxviii}{214339}

\newcommand{\WallLjxxxviii}{2.204604}

\newcommand{\IqrWallLjxxxviii}{0.175856}

\newcommand{\HitLjlv}{6}

\newcommand{\MedLjlv}{5828.5}

\newcommand{\IqrLjlv}{7776}

\newcommand{\WallLjlv}{3.833091}

\newcommand{\IqrWallLjlv}{0.364198}

\newcommand{\HitLjlxxv}{1}

\newcommand{\MedLjlxxv}{219506}

\newcommand{\IqrLjlxxv}{0}

\newcommand{\WallLjlxxv}{9.148674}

\newcommand{\IqrWallLjlxxv}{2.127980}

\newcommand{\HitLjxcviii}{0}

\newcommand{\MedLjxcviii}{219506}

\newcommand{\IqrLjxcviii}{0}

\newcommand{\WallLjxcviii}{13.252215}

\newcommand{\IqrWallLjxcviii}{2.884105}

\newcommand{\HitLjcv}{0}

\newcommand{\MedLjcv}{219506}

\newcommand{\IqrLjcv}{0}

\newcommand{\WallLjcv}{18.297832}

\newcommand{\IqrWallLjcv}{2.195455}

\title{Supporting Information: Recommended basin hopping on Lennard-Jones clusters}
\author{Rohit Goswami}
% vale rgoswami.AbbreviationStyle = NO
\affiliation{SURF B.V., Amsterdam, The Netherlands}
% vale rgoswami.AbbreviationStyle = YES
\alsoaffiliation{Science Institute, University of Iceland, Reykjavik, Iceland}
\alsoaffiliation{TurtleTech ehf., Reykjavik, Iceland}
\author{Ruhila Goswami}
\affiliation{Faculty of Life and Environmental Sciences, University of Iceland, Reykjavik, Iceland}
\author{Amrita Goswami}
\affiliation{Universidad Complutense de Madrid, Madrid, Spain}
\author{Moritz Sallermann}
\affiliation{Universidad Complutense de Madrid, Madrid, Spain}
\author{Sonaly Goswami}
\affiliation{Department of Chemistry, Indian Institute of Technology Kanpur, Kanpur, India}
\author{Debabrata Goswami}
\affiliation{Center for Bits and Atoms, Massachusetts Institute of Technology, Cambridge, Massachusetts 02139, United States}
\alsoaffiliation{Department of Chemistry, Indian Institute of Technology Kanpur, Kanpur, India}
\email{dgoswami@iitk.ac.in}
\begin{document}

\section{Board constants}
The hit rule is an energy strictly below the reference plus \(0.001\), read from an improvement in the trace.
The references are \(-173.928427\) (LJ\(_{38}\)), \(-279.248470\) (LJ\(_{55}\)), \(-397.492331\) (LJ\(_{75}\)), \(-543.665361\) (LJ\(_{98}\)) and \(-582.086642\) (LJ\(_{104}\)), in reduced units.
The temperature is \(0.8\), the value Wales and Doye used for Lennard-Jones basin hopping and the value stored at \texttt{src/methods/cluster\_hopping/preset.rs:8}.\cite{wales1997}
The record gradient is \(0.001\), stored at \texttt{src/methods/cluster\_hopping/preset.rs:31}.
The recommended quench length is 1000 steps, assigned at \texttt{src/methods/cluster\_hopping/config.rs:1138}.
\texttt{Config::for\_cluster} keeps 200 steps, assigned at \texttt{src/methods/cluster\_hopping/config.rs:1425}.
The start density passed to \texttt{random\_cluster} is \(0.7\), the constant \texttt{LJ\_LIQUID\_NUMBER\_DENSITY} at \texttt{src/methods/cluster\_hopping.rs:6502}.
Thol and coworkers give 0.702 from their ancillary equation at reduced temperature 1, and the container keeps one decimal place.
The call ceiling is \CallCeiling{} at every size.
For six sorted values the median is the mean of the third and fourth, and the interquartile range is the fifth minus the second.
A miss is charged twice its ceiling, which is \MissPar{}.

\section{Six fresh runs}
The binary was built from commit \texttt{\CommitSha}.
That commit drops the atom-count branches from \texttt{Config::recommended} and from \texttt{Config::communicating}, and it names the 1000-step quench.
Main, commit \texttt{\MainCommit}, still installs the reseed library at 38 atoms and a compacted surface below 98 atoms.
Every log header records \texttt{uname -m}, \texttt{nproc}, and the processor string.
The launch order is the file \texttt{logs/order.txt}.
The library on every run is \texttt{LeanBurst}, and the surface count is 0.
At 38 atoms the sorted penalised costs are 2421, 5167, 26528, 51767, 219506 and 219506.
The median is \((26528+51767)/2 = 39147.5\).
The interquartile range is \(219506-5167 = 214339\).
\begin{center}
\begin{tabular}{rrrrrr}
\toprule
size & seed & hit & calls & PAR-2 & wall / s \\
\midrule
38 & 0 & 1 & 2421 & 2421 & 2.049692 \\
38 & 1 & 1 & 5167 & 5167 & 2.215312 \\
38 & 2 & 0 & 109753 & 219506 & 2.284813 \\
38 & 3 & 0 & 109753 & 219506 & 2.321189 \\
38 & 4 & 1 & 51767 & 51767 & 2.108957 \\
38 & 5 & 1 & 26528 & 26528 & 2.193896 \\
55 & 0 & 1 & 6251 & 6251 & 3.878319 \\
55 & 1 & 1 & 13742 & 13742 & 3.787863 \\
55 & 2 & 1 & 2823 & 2823 & 3.553142 \\
55 & 3 & 1 & 11288 & 11288 & 4.242409 \\
55 & 4 & 1 & 5406 & 5406 & 3.634336 \\
55 & 5 & 1 & 3512 & 3512 & 3.998534 \\
75 & 0 & 0 & 109753 & 219506 & 10.356952 \\
75 & 1 & 0 & 109753 & 219506 & 8.196650 \\
75 & 2 & 0 & 109753 & 219506 & 8.134675 \\
75 & 3 & 0 & 109753 & 219506 & 9.454797 \\
75 & 4 & 0 & 109753 & 219506 & 8.842551 \\
75 & 5 & 1 & 60314 & 60314 & 10.324630 \\
98 & 0 & 0 & 109753 & 219506 & 12.278025 \\
98 & 1 & 0 & 109753 & 219506 & 13.052619 \\
98 & 2 & 0 & 109753 & 219506 & 15.368852 \\
98 & 3 & 0 & 109753 & 219506 & 13.451811 \\
98 & 4 & 0 & 109753 & 219506 & 12.352151 \\
98 & 5 & 0 & 109753 & 219506 & 15.236256 \\
104 & 0 & 0 & 109753 & 219506 & 18.211995 \\
104 & 1 & 0 & 109753 & 219506 & 19.429353 \\
104 & 2 & 0 & 109753 & 219506 & 13.449188 \\
104 & 3 & 0 & 109753 & 219506 & 17.233898 \\
104 & 4 & 0 & 109753 & 219506 & 22.499415 \\
104 & 5 & 0 & 109753 & 219506 & 18.383669 \\
\bottomrule
\end{tabular}

\end{center}

\section{Published symmetrised costs}
\texttt{data/literature.tsv} quotes Oakley, Johnston and Wales, Table 1.\cite{oakley2013}
The continuous-symmetry basin-hopping mean at 38 atoms is 4369 evaluations over 100 successful runs.
The core-orbit basin-hopping means are 50229 at 75 atoms and 109753 at 98 atoms, again over 100 successful runs.
The same file's Table 2 rows are genetic-algorithm means at 75 and 98 atoms.
No row in that file is a genetic algorithm at 104 atoms.

\section{TIP4P pair constants}
\texttt{logs/tip4p\_constants.txt} records Jorgensen's \(A = 2510400\) and \(C = 2552.24\), in kJ\,\AA\(^{12}\)\,mol\(^{-1}\) and kJ\,\AA\(^{6}\)\,mol\(^{-1}\).\cite{jorgensen1983}
The global minimisation program matches those two numbers for \(C_6\) and \(C_{12}\).
The anneal pair built from a rounded well depth and a rounded \(\sigma\) gives \(C_6 = 2551.8252\) and \(C_{12} = 2510336.2\).
This article reports no TIP4P hit count.
A later TIP4P campaign has to use one pair of coefficients on every code.

\section{Rerun plan}
We have not run the eight experiments in this list.
E1 defines the calls-to-target of a multi-replica run, then runs one chain, four independent chains, four communicating chains, and a bank of forty chains, on the shared starts of LJ\(_{38}\), LJ\(_{55}\), LJ\(_{75}\), LJ\(_{98}\) and TIP4P\(_{10}\).
E2 runs the global minimisation program with the archived keyword file for each size, dynamic lattice search, a cut-and-splice genetic algorithm, ABCluster, and minima hopping, on the same starts, with one force-call ledger.
E3 adds an arm with retirement off and an arm that is plain basin hopping plus orbit completion, and scores them where a screen can refuse a trial before the hit.
E4 uses one hit rule on every code, a minimum relaxed to a stated gradient bound, and pairs the archived keyword files on the shared starts.
E5 is the constructor measured in the article, with the 1000-step quench and the record gradient 0.001, scored again on held-out sizes from LJ\(_{31}\) to LJ\(_{37}\), LJ\(_{76}\) to LJ\(_{97}\), and LJ\(_{102}\) to LJ\(_{104}\).
E6 uses Jorgensen's \(A\) and \(C\) on every TIP4P code and scores every size at \(0.001\,\mathrm{kJ\,mol^{-1}}\).
E7 runs the external potential path and the built-in surface on one system and compares hits and calls.
E8 is a bond-constrained conformer comparison against CREST on valid records, or the absence of that section.
Hit counts in that campaign are compared by Fisher's exact test.
Paired penalised costs are compared by the Wilcoxon signed-rank test, with Holm's correction over a family written down before the run.\cite{holm1979}

\section{Repeating the six runs}
Build \texttt{examples/lj38\_honest} in release and run, for each pair in \texttt{logs/order.txt},
\begin{verbatim}
SEED_OFFSET=SEED RAYON_NUM_THREADS=1 taskset -c 0 lj38_honest N SEED 109753
\end{verbatim}
The ceiling is \CallCeiling{} at every size.
A header with \texttt{uname -m}, \texttt{nproc}, and the processor string is written on every log, and the hostname is left out.

\bibliography{refs}

\end{document}
